% GENERATED FILE. Edit canonical lessons, then run npm run generate:books.
% canonical-source-hash: fnv1a64:4a8cf35055409110
% canonical-lessons: ZH-C500-huosheng, ZH-C500-qiusai, ZH-W500-sai-match, ZH-C500-juesai, ZH-C500-bifen, ZH-C500-shangchang

\chapter{Match Day}
\label{ch:zh-match-day}
\hlchaptermodality{500}

\begin{chapteropening}
\textbf{By the end of this chapter:} \emph{I can talk about winning, a match, a final, the score and taking the field.}

Complete the last lesson of chapter 500: I can talk about winning, a match, a final, the score and taking the field.
\end{chapteropening}

\section[huòshèng]{huòshèng — to win a contest; to be victorious}

\label{lesson:ZH-C500-huosheng}



\begin{quote}
Before the new one: say the Mandarin for to do training (for a sport), then the Mandarin for to pass the ball.
\end{quote}

\subsection*{You'll want to know: huòshèng}
\emph{huòshèng} — \textquotedblleft{}to win a contest; to be victorious\textquotedblright{}.

\begin{quote}
\emph{wǒmen de qiúduì huòshèng le wǒ gǎnjué tèbié gāoxìng} — Our team won, and I'm over the moon.
\end{quote}

Its characters come later; for now it is a word for the ear.

\begin{cousinweb}
\emph{huòshèng} is built from \emph{huò} 'to get' and \emph{shèng} 'victory', both inherited from Middle Chinese, \emph{hweak} and \emph{syingH} (Guangyun, in Baxter's transcription), which glosses the first 'to get' and the second 'victory and defeat' and 'to overcome'.
\end{cousinweb}

\begin{grammarlens}[title={What you've built}]
One more word for saying how it was.
\end{grammarlens}

\subsection*{Your turn}
\begin{itemize}
\raggedright
  \item \emph{Say it:} \emph{huòshèng}
  \item \emph{Say it:} \emph{huòshèng}, once more
  \item \emph{Say it:} \emph{wǒmen de qiúduì huòshèng le wǒ gǎnjué tèbié gāoxìng} — Our team won, and I'm over the moon.
  \item \emph{Recall:} say the Mandarin for to do training (for a sport), then the Mandarin for to pass the ball, then say \emph{huòshèng} again
\end{itemize}

\begin{tcolorbox}[breakable,colback=teal!4,colframe=teal!35!black,title={Before you move on}]
What does huòshèng mean? (\textquotedblleft{}To win a contest; to be victorious\textquotedblright{}.) Say it once more.
\end{tcolorbox}
\section[\texorpdfstring{\zh{球赛} (qiúsài)}{球赛 (qiúsài)}]{\zh{球赛} (qiúsài) — a ball game (a football or basketball match)}

\label{lesson:ZH-C500-qiusai}



\begin{quote}
Before the new one: say the Mandarin for to pass the ball, then the Mandarin for to win a contest.
\end{quote}

\subsection*{You'll want to know: \zh{球赛}}
\textbf{\zh{球赛}} — \emph{qiúsài} — \textquotedblleft{}a ball game (a football or basketball match)\textquotedblright{}.

\begin{quote}
\textbf{\zh{这是我看过的最好的一场球赛}} \emph{zhè shì wǒ kàn guo de zuì hǎo de yì chǎng qiúsài} — This is the best match I've ever seen.
\end{quote}

The new character, in pieces: \zh{赛} — fourteen strokes in three pieces: the roof of \zh{家} \emph{jiā} on top, three strokes; in the middle seven strokes, two short bars, two verticals, a long bar and two strokes spreading left and right; at the foot four strokes. The next lesson puts it in your hand.

\begin{cousinweb}
\textbf{\zh{球赛}} (\emph{qiúsài}) is built from \zh{球} 'a ball' and \zh{赛} 'a contest', both inherited from Middle Chinese, \emph{gjuw} and \emph{sojH} (Guangyun, in Baxter's transcription), which glosses the second 'to repay'; the sense 'a contest' is not given there.
\end{cousinweb}

\begin{grammarlens}[title={What you've built}]
One more word for saying how it was.
\end{grammarlens}

\subsection*{Your turn}
\begin{itemize}
\raggedright
  \item \emph{Say it:} \emph{qiúsài}
  \item \emph{Say it:} \emph{qiúsài}, once more
  \item \emph{Say it:} \emph{zhè shì wǒ kàn guo de zuì hǎo de yì chǎng qiúsài} — This is the best match I've ever seen.
  \item \emph{Recall:} say the Mandarin for to pass the ball, then the Mandarin for to win a contest, then say \emph{qiúsài} again
\end{itemize}

\begin{tcolorbox}[breakable,colback=teal!4,colframe=teal!35!black,title={Before you move on}]
What does \zh{球赛} mean? (\textquotedblleft{}A ball game (a football or basketball match)\textquotedblright{}.) Say it once more.
\end{tcolorbox}
\section[\texorpdfstring{\zh{赛} (sài)}{赛 (sài)}]{\zh{赛} — a roof, a middle of seven strokes, a foot of four}

\label{lesson:ZH-W500-sai-match}



\begin{quote}
The word first, then the character inside it.

\begin{itemize}
\raggedright
  \item \emph{Recall:} say \emph{qiúsài} — \textbf{R1}, one lesson back, before it decays
\end{itemize}

\begin{quote}
\zh{家} \zh{比}
\end{quote}

\emph{Read it:} each one aloud, with nothing said first — \emph{jiā}, \emph{bǐ}
\end{quote}

\subsection*{Script — a roof, a middle of seven strokes, a foot of four}
\hlblockfigure{\detokenize{figures/ZH-W500-sai-match-filmstrip.pdf}}{How \zh{赛} is written, stroke by stroke}

\begin{quote}
\zh{赛}
\end{quote}

\textbf{Three pieces, top to bottom.}

\textbf{The roof, three strokes}, as on top of \zh{家} and \zh{安}: a dot, a short stroke down on the left, then a bar that turns down at the right end and hooks.

\textbf{The middle, seven strokes.} Two short bars, one under the other. Two verticals down through them. A long bar across. Then a stroke falling to the left and a stroke falling to the right, spreading out like legs.

\textbf{The foot, four strokes}: a vertical, a bar that turns down, a stroke falling to the left inside, and a dot.

3 + 7 + 4 = fourteen.

\subsection*{Your turn}
\begin{itemize}
\raggedright
  \item \emph{Write it:} \zh{赛}, all fourteen, with the model in front of you
  \item \emph{Cover:} write it again from memory, and count the pieces, not the strokes
\end{itemize}

\begin{tcolorbox}[breakable,colback=teal!4,colframe=teal!35!black,title={Before you move on}]
How many pieces does \zh{赛} have? (\textbf{Three}.) What is on top? (\textbf{The roof of \zh{家}}.) How many strokes at the foot? (\textbf{Four}.)

Source: \href{https://raw.githubusercontent.com/chanind/hanzi-writer-data/68d10a4b21150cae5e1ebbd223eed289cf32d90c/data/\%E8\%B5\%9B.json}{Hanzi Writer Data \zh{赛}, pinned stroke paths and medians}.
\end{tcolorbox}
\section[juésài]{juésài — final (of a competition)}

\label{lesson:ZH-C500-juesai}



\begin{quote}
Before the new one: say the Mandarin for to win a contest, then the Mandarin for a ball game.
\end{quote}

\subsection*{You'll want to know: juésài}
\emph{juésài} — \textquotedblleft{}final (of a competition)\textquotedblright{}.

\begin{quote}
\emph{wǒmen de qiúduì jìn le juésài} — Our team got into the final.
\end{quote}

Its characters come later; for now it is a word for the ear.

\begin{cousinweb}
\emph{juésài} is built from \emph{jué} 'to decide' and \zh{赛} 'a contest', both inherited from Middle Chinese, \emph{kwet} and \emph{sojH} (Guangyun, in Baxter's transcription), which glosses the first 'to cut' and 'to break through'.
\end{cousinweb}

\begin{grammarlens}[title={What you've built}]
One more word for saying how it was.
\end{grammarlens}

\subsection*{Your turn}
\begin{itemize}
\raggedright
  \item \emph{Say it:} \emph{juésài}
  \item \emph{Say it:} \emph{juésài}, once more
  \item \emph{Say it:} \emph{wǒmen de qiúduì jìn le juésài} — Our team got into the final.
  \item \emph{Recall:} say the Mandarin for to win a contest, then the Mandarin for a ball game, then say \emph{juésài} again
\end{itemize}

\begin{tcolorbox}[breakable,colback=teal!4,colframe=teal!35!black,title={Before you move on}]
What does juésài mean? (\textquotedblleft{}Final (of a competition)\textquotedblright{}.) Say it once more.
\end{tcolorbox}
\section[bǐfēn]{bǐfēn — the score of a game}

\label{lesson:ZH-C500-bifen}



\begin{quote}
Before the new one: say the Mandarin for a ball game, then the Mandarin for final (of a competition).
\end{quote}

\subsection*{You'll want to know: bǐfēn}
\emph{bǐfēn} — \textquotedblleft{}the score of a game\textquotedblright{}.

\begin{quote}
\emph{xiànzài bǐfēn shì duōshao} — What's the score now?
\end{quote}

Its characters come later; for now it is a word for the ear.

\begin{cousinweb}
\emph{bǐfēn} is built from \zh{比} 'to compare' and \emph{fēn} 'a part', both inherited from Middle Chinese, \emph{pjijX} and \emph{pjun} (Guangyun, in Baxter's transcription), which glosses the first 'to collate' and the second 'to divide'.
\end{cousinweb}

\begin{grammarlens}[title={What you've built}]
One more word for saying how it was.
\end{grammarlens}

\subsection*{Your turn}
\begin{itemize}
\raggedright
  \item \emph{Say it:} \emph{bǐfēn}
  \item \emph{Say it:} \emph{bǐfēn}, once more
  \item \emph{Say it:} \emph{xiànzài bǐfēn shì duōshao} — What's the score now?
  \item \emph{Recall:} say the Mandarin for a ball game, then the Mandarin for final (of a competition), then say \emph{bǐfēn} again
\end{itemize}

\begin{tcolorbox}[breakable,colback=teal!4,colframe=teal!35!black,title={Before you move on}]
What does bǐfēn mean? (\textquotedblleft{}The score of a game\textquotedblright{}.) Say it once more.
\end{tcolorbox}
\section[\texorpdfstring{\zh{上场} (shàngchǎng)}{上场 (shàngchǎng)}]{\zh{上场} (shàngchǎng) — to take the field; to come on (to play)}

\label{lesson:ZH-C500-shangchang}



\begin{quote}
Before the new one: say the Mandarin for final (of a competition), then the Mandarin for the score of a game.
\end{quote}

\subsection*{You'll want to know: \zh{上场}}
\textbf{\zh{上场}} — \emph{shàngchǎng} — \textquotedblleft{}to take the field; to come on (to play)\textquotedblright{}.

\begin{quote}
\textbf{\zh{这场比赛我没有上场}} \emph{zhè chǎng bǐsài wǒ méiyǒu shàngchǎng} — I didn't get on the pitch in this match.
\end{quote}

\begin{cousinweb}
\textbf{\zh{上场}} (\emph{shàngchǎng}) is built from \zh{上} 'to go up' and \zh{场} 'a ground', both inherited from Middle Chinese, \emph{dzyangX} and \emph{drjang} (Guangyun, in Baxter's transcription), which glosses the first 'to climb' and the second 'a threshing floor'.
\end{cousinweb}

\begin{grammarlens}[title={What you've built}]
One more word for saying how it was.
\end{grammarlens}

\subsection*{Your turn}
\begin{itemize}
\raggedright
  \item \emph{Say it:} \emph{shàngchǎng}
  \item \emph{Say it:} \emph{shàngchǎng}, once more
  \item \emph{Say it:} \emph{zhè chǎng bǐsài wǒ méiyǒu shàngchǎng} — I didn't get on the pitch in this match.
  \item \emph{Recall:} say the Mandarin for final (of a competition), then the Mandarin for the score of a game, then say \emph{shàngchǎng} again
\end{itemize}

\begin{tcolorbox}[breakable,colback=teal!4,colframe=teal!35!black,title={Before you move on}]
What does \zh{上场} mean? (\textquotedblleft{}To take the field; to come on (to play)\textquotedblright{}.) Say it once more.
\end{tcolorbox}
